\subsection{Controlled-Replay Negative Result and Capability Boundary}

The replay experiment distinguishes conditional ranking ability from end-to-end diagnostic ability. When the known injection onset and a matched baseline were supplied, the topic-to-module evidence placed the injected module within the Top-1 for some EKF2 runs, but performance was uneven across mutation types and the aggregate Top-5 recall remained 0.333. This onset-conditioned analysis tests whether a localized change in telemetry evidence can be propagated toward a relevant module after temporal alignment has already been provided. It does not test whether the deployed chain can independently recognize the change.

The detector-gated result provides that stronger test and constitutes the main capability boundary of the current system. The unchanged real-flight detector triggered on none of the 12 mutation runs, so all detector-gated ranking metrics were zero. This failure cannot be offset by the onset-conditioned ranks because an online chain cannot produce a timely module suspect when its first stage does not identify the input as anomalous. Reporting the two evaluation layers separately prevents privileged onset information from being mistaken for operational localization performance.

A plausible interpretation is distribution mismatch between the real-flight annotations used to train Stage 1 and the controlled source mutations used in replay. The two settings differ in anomaly mechanisms, signal manifestations, and replay conditions, and the replay benchmark was not used to adapt the detector. This interpretation remains a hypothesis rather than a quantified decomposition of domain shift. The negative result nevertheless establishes that evidence propagation alone is insufficient: end-to-end module diagnosis requires a detector whose operating distribution includes, or generalizes to, the relevant software-induced behaviors.
